% !TeX spellcheck = de_DE_frami
%%%%%%%%%%%%%%%%%%%%%%%%%%%%%%%%%%%%%%%%%%%%%%%%
% COPYRIGHT: (C) 2017-2026 FAU FabLab and others
% CC-BY-SA 3.0
% Text in eigenen Worten, KEINE Passagen oder Abbildungen aus der
% Proxxon-Betriebsanleitung übernehmen (siehe README, Abschnitt Lizenz).
%%%%%%%%%%%%%%%%%%%%%%%%%%%%%%%%%%%%%%%%%%%%%%%%

\newcommand{\basedir}{fablab-document}
\documentclass{\basedir/fablab-document}

\usepackage{amssymb}
\usepackage{textcomp}
\usepackage{tabularx} % Tabellen mit bestimmtem Breitenverhältnis der Spalten
\usepackage{float}
\usepackage{caption}
\usepackage{qrcode}
\usepackage{fablab-document/fablab-ba} % Betriebsanweisung
% =====================================================================
%  Gemeinsame Farben und Stile der Zeichnungen (TikZ, selbst erstellt)
%  In der Präambel einbinden: % =====================================================================
%  Gemeinsame Farben und Stile der Zeichnungen (TikZ, selbst erstellt)
%  In der Präambel einbinden: % =====================================================================
%  Gemeinsame Farben und Stile der Zeichnungen (TikZ, selbst erstellt)
%  In der Präambel einbinden: \input{zeichnungen/stil}
%  Angelehnt an die Zeichnungen der anderen FabLab-Einweisungen
%  (z.B. fau-fablab/shaper-origin-einweisung). Lizenz: CC BY-SA 3.0
% =====================================================================
\usetikzlibrary{calc,arrows.meta,decorations.pathreplacing,patterns}

\definecolor{okgruen}{RGB}{46,160,67}
\definecolor{nokrot}{RGB}{207,34,46}
\definecolor{mittel}{RGB}{230,150,0}
\definecolor{feld}{RGB}{244,246,248}
\definecolor{holz}{RGB}{226,196,150}
\definecolor{holzrand}{RGB}{150,110,60}
\definecolor{gehaeusegruen}{RGB}{88,110,80}   % Gehäuse der Säge
\definecolor{gelbteil}{RGB}{226,200,40}       % gelbe Bedienteile
\definecolor{schutzorange}{RGB}{240,120,40}   % Sägeblattschutz
\definecolor{alu}{RGB}{205,210,215}           % Sägetisch
\definecolor{gefahr}{RGB}{207,34,46}

\tikzset{
  % Bauteile der Maschine
  gehaeuse/.style={fill=gehaeusegruen!70, draw=black!80, line width=0.7pt, line join=round},
  tisch/.style={fill=alu, draw=black!75, line width=0.7pt, line join=round},
  metall/.style={fill=black!30, draw=black!75, line width=0.6pt, line join=round},
  gelb/.style={fill=gelbteil, draw=black!75, line width=0.6pt, line join=round},
  blatt/.style={fill=black!55, draw=black!85, line width=0.5pt, line join=round},
  keil/.style={fill=black!20, draw=black!85, line width=0.6pt, line join=round},
  werkstueck/.style={fill=holz, draw=holzrand, line width=0.6pt, line join=round},
  hand/.style={fill=orange!20, draw=orange!60!black, line width=0.6pt, line join=round, rounded corners=1.2mm},
  gefahrenzone/.style={pattern=north east lines, pattern color=gefahr!70, draw=gefahr, dashed, line width=0.5pt},
  % Pfeile, Maße, Beschriftung
  vorschub/.style={-{Stealth[length=2.6mm, width=2.2mm]}, line width=1.3pt},
  drehung/.style={-{Stealth[length=1.8mm, width=1.6mm]}, line width=0.9pt},
  kraft/.style={-{Stealth[length=2.2mm, width=2mm]}, line width=1.1pt},
  mass/.style={{Stealth[length=1.4mm]}-{Stealth[length=1.4mm]}, line width=0.4pt, black!80},
  hilfslinie/.style={line width=0.3pt, black!60},
  marke/.style={circle, draw=black, fill=white, line width=0.8pt,
    minimum size=4.6mm, inner sep=0pt, font=\scriptsize\bfseries},
  zeiger/.style={line width=0.6pt},
  % Bewertung: \pic at (x,y) {ok}; bzw. {nok}
  ok/.pic={\fill[okgruen] (0,0) circle (3);
    \draw[white, line width=1.3pt, line cap=round, line join=round]
      (-1.4,0.1) -- (-0.4,-1.1) -- (1.5,1.2);},
  nok/.pic={\fill[nokrot] (0,0) circle (3);
    \draw[white, line width=1.3pt, line cap=round]
      (-1.1,-1.1) -- (1.1,1.1) (-1.1,1.1) -- (1.1,-1.1);},
}

% Feld mit hellem Hintergrund und Überschrift: \feld{x}{y}{Breite}{Höhe}{Titel}
\newcommand{\feld}[5]{%
  \fill[feld, rounded corners=2mm] (#1,#2) rectangle ++(#3,#4);
  \node[anchor=north west, font=\small\bfseries] at (#1+3,#2+#4-2) {#5};}

% Nummer an einer Zeichnung mit Zeigerlinie: \nr{Nummer}{Position}{Ziel}
\newcommand{\nr}[3]{%
  \draw[zeiger] (#2) -- (#3);
  \fill (#3) circle[radius=0.6];
  \node[marke] at (#2) {#1};}

% Kreissägeblatt in der Seitenansicht: \saegeblatt{x}{y}{Radius}{Zähne}
% Die Zähne zeigen im Uhrzeigersinn. In allen Seitenansichten steht die
% Bedienperson rechts und schiebt das Werkstück nach links: Die Zähne
% laufen oben auf sie zu und schneiden vorn von oben nach unten.
\newcommand{\saegeblatt}[4]{%
  \begin{scope}[shift={(#1,#2)}]
    \pgfmathsetmacro{\sbs}{360/#4}%
    \pgfmathsetmacro{\sbz}{min(0.18*#3, 2*pi*#3/#4*0.55)}%
    \fill[black!60, draw=black!85, line width=0.4pt] (0:#3)
      \foreach \i in {1,...,#4}{
        -- ({\i*\sbs-\sbs}:#3) -- ({\i*\sbs-0.2*\sbs}:#3-\sbz) -- ({\i*\sbs}:#3)}
      -- cycle;
    \fill[black!25] (0,0) circle (#3*0.5);
    \fill[white] (0,0) circle (#3*0.12);
  \end{scope}}

%  Angelehnt an die Zeichnungen der anderen FabLab-Einweisungen
%  (z.B. fau-fablab/shaper-origin-einweisung). Lizenz: CC BY-SA 3.0
% =====================================================================
\usetikzlibrary{calc,arrows.meta,decorations.pathreplacing,patterns}

\definecolor{okgruen}{RGB}{46,160,67}
\definecolor{nokrot}{RGB}{207,34,46}
\definecolor{mittel}{RGB}{230,150,0}
\definecolor{feld}{RGB}{244,246,248}
\definecolor{holz}{RGB}{226,196,150}
\definecolor{holzrand}{RGB}{150,110,60}
\definecolor{gehaeusegruen}{RGB}{88,110,80}   % Gehäuse der Säge
\definecolor{gelbteil}{RGB}{226,200,40}       % gelbe Bedienteile
\definecolor{schutzorange}{RGB}{240,120,40}   % Sägeblattschutz
\definecolor{alu}{RGB}{205,210,215}           % Sägetisch
\definecolor{gefahr}{RGB}{207,34,46}

\tikzset{
  % Bauteile der Maschine
  gehaeuse/.style={fill=gehaeusegruen!70, draw=black!80, line width=0.7pt, line join=round},
  tisch/.style={fill=alu, draw=black!75, line width=0.7pt, line join=round},
  metall/.style={fill=black!30, draw=black!75, line width=0.6pt, line join=round},
  gelb/.style={fill=gelbteil, draw=black!75, line width=0.6pt, line join=round},
  blatt/.style={fill=black!55, draw=black!85, line width=0.5pt, line join=round},
  keil/.style={fill=black!20, draw=black!85, line width=0.6pt, line join=round},
  werkstueck/.style={fill=holz, draw=holzrand, line width=0.6pt, line join=round},
  hand/.style={fill=orange!20, draw=orange!60!black, line width=0.6pt, line join=round, rounded corners=1.2mm},
  gefahrenzone/.style={pattern=north east lines, pattern color=gefahr!70, draw=gefahr, dashed, line width=0.5pt},
  % Pfeile, Maße, Beschriftung
  vorschub/.style={-{Stealth[length=2.6mm, width=2.2mm]}, line width=1.3pt},
  drehung/.style={-{Stealth[length=1.8mm, width=1.6mm]}, line width=0.9pt},
  kraft/.style={-{Stealth[length=2.2mm, width=2mm]}, line width=1.1pt},
  mass/.style={{Stealth[length=1.4mm]}-{Stealth[length=1.4mm]}, line width=0.4pt, black!80},
  hilfslinie/.style={line width=0.3pt, black!60},
  marke/.style={circle, draw=black, fill=white, line width=0.8pt,
    minimum size=4.6mm, inner sep=0pt, font=\scriptsize\bfseries},
  zeiger/.style={line width=0.6pt},
  % Bewertung: \pic at (x,y) {ok}; bzw. {nok}
  ok/.pic={\fill[okgruen] (0,0) circle (3);
    \draw[white, line width=1.3pt, line cap=round, line join=round]
      (-1.4,0.1) -- (-0.4,-1.1) -- (1.5,1.2);},
  nok/.pic={\fill[nokrot] (0,0) circle (3);
    \draw[white, line width=1.3pt, line cap=round]
      (-1.1,-1.1) -- (1.1,1.1) (-1.1,1.1) -- (1.1,-1.1);},
}

% Feld mit hellem Hintergrund und Überschrift: \feld{x}{y}{Breite}{Höhe}{Titel}
\newcommand{\feld}[5]{%
  \fill[feld, rounded corners=2mm] (#1,#2) rectangle ++(#3,#4);
  \node[anchor=north west, font=\small\bfseries] at (#1+3,#2+#4-2) {#5};}

% Nummer an einer Zeichnung mit Zeigerlinie: \nr{Nummer}{Position}{Ziel}
\newcommand{\nr}[3]{%
  \draw[zeiger] (#2) -- (#3);
  \fill (#3) circle[radius=0.6];
  \node[marke] at (#2) {#1};}

% Kreissägeblatt in der Seitenansicht: \saegeblatt{x}{y}{Radius}{Zähne}
% Die Zähne zeigen im Uhrzeigersinn. In allen Seitenansichten steht die
% Bedienperson rechts und schiebt das Werkstück nach links: Die Zähne
% laufen oben auf sie zu und schneiden vorn von oben nach unten.
\newcommand{\saegeblatt}[4]{%
  \begin{scope}[shift={(#1,#2)}]
    \pgfmathsetmacro{\sbs}{360/#4}%
    \pgfmathsetmacro{\sbz}{min(0.18*#3, 2*pi*#3/#4*0.55)}%
    \fill[black!60, draw=black!85, line width=0.4pt] (0:#3)
      \foreach \i in {1,...,#4}{
        -- ({\i*\sbs-\sbs}:#3) -- ({\i*\sbs-0.2*\sbs}:#3-\sbz) -- ({\i*\sbs}:#3)}
      -- cycle;
    \fill[black!25] (0,0) circle (#3*0.5);
    \fill[white] (0,0) circle (#3*0.12);
  \end{scope}}

%  Angelehnt an die Zeichnungen der anderen FabLab-Einweisungen
%  (z.B. fau-fablab/shaper-origin-einweisung). Lizenz: CC BY-SA 3.0
% =====================================================================
\usetikzlibrary{calc,arrows.meta,decorations.pathreplacing,patterns}

\definecolor{okgruen}{RGB}{46,160,67}
\definecolor{nokrot}{RGB}{207,34,46}
\definecolor{mittel}{RGB}{230,150,0}
\definecolor{feld}{RGB}{244,246,248}
\definecolor{holz}{RGB}{226,196,150}
\definecolor{holzrand}{RGB}{150,110,60}
\definecolor{gehaeusegruen}{RGB}{88,110,80}   % Gehäuse der Säge
\definecolor{gelbteil}{RGB}{226,200,40}       % gelbe Bedienteile
\definecolor{schutzorange}{RGB}{240,120,40}   % Sägeblattschutz
\definecolor{alu}{RGB}{205,210,215}           % Sägetisch
\definecolor{gefahr}{RGB}{207,34,46}

\tikzset{
  % Bauteile der Maschine
  gehaeuse/.style={fill=gehaeusegruen!70, draw=black!80, line width=0.7pt, line join=round},
  tisch/.style={fill=alu, draw=black!75, line width=0.7pt, line join=round},
  metall/.style={fill=black!30, draw=black!75, line width=0.6pt, line join=round},
  gelb/.style={fill=gelbteil, draw=black!75, line width=0.6pt, line join=round},
  blatt/.style={fill=black!55, draw=black!85, line width=0.5pt, line join=round},
  keil/.style={fill=black!20, draw=black!85, line width=0.6pt, line join=round},
  werkstueck/.style={fill=holz, draw=holzrand, line width=0.6pt, line join=round},
  hand/.style={fill=orange!20, draw=orange!60!black, line width=0.6pt, line join=round, rounded corners=1.2mm},
  gefahrenzone/.style={pattern=north east lines, pattern color=gefahr!70, draw=gefahr, dashed, line width=0.5pt},
  % Pfeile, Maße, Beschriftung
  vorschub/.style={-{Stealth[length=2.6mm, width=2.2mm]}, line width=1.3pt},
  drehung/.style={-{Stealth[length=1.8mm, width=1.6mm]}, line width=0.9pt},
  kraft/.style={-{Stealth[length=2.2mm, width=2mm]}, line width=1.1pt},
  mass/.style={{Stealth[length=1.4mm]}-{Stealth[length=1.4mm]}, line width=0.4pt, black!80},
  hilfslinie/.style={line width=0.3pt, black!60},
  marke/.style={circle, draw=black, fill=white, line width=0.8pt,
    minimum size=4.6mm, inner sep=0pt, font=\scriptsize\bfseries},
  zeiger/.style={line width=0.6pt},
  % Bewertung: \pic at (x,y) {ok}; bzw. {nok}
  ok/.pic={\fill[okgruen] (0,0) circle (3);
    \draw[white, line width=1.3pt, line cap=round, line join=round]
      (-1.4,0.1) -- (-0.4,-1.1) -- (1.5,1.2);},
  nok/.pic={\fill[nokrot] (0,0) circle (3);
    \draw[white, line width=1.3pt, line cap=round]
      (-1.1,-1.1) -- (1.1,1.1) (-1.1,1.1) -- (1.1,-1.1);},
}

% Feld mit hellem Hintergrund und Überschrift: \feld{x}{y}{Breite}{Höhe}{Titel}
\newcommand{\feld}[5]{%
  \fill[feld, rounded corners=2mm] (#1,#2) rectangle ++(#3,#4);
  \node[anchor=north west, font=\small\bfseries] at (#1+3,#2+#4-2) {#5};}

% Nummer an einer Zeichnung mit Zeigerlinie: \nr{Nummer}{Position}{Ziel}
\newcommand{\nr}[3]{%
  \draw[zeiger] (#2) -- (#3);
  \fill (#3) circle[radius=0.6];
  \node[marke] at (#2) {#1};}

% Kreissägeblatt in der Seitenansicht: \saegeblatt{x}{y}{Radius}{Zähne}
% Die Zähne zeigen im Uhrzeigersinn. In allen Seitenansichten steht die
% Bedienperson rechts und schiebt das Werkstück nach links: Die Zähne
% laufen oben auf sie zu und schneiden vorn von oben nach unten.
\newcommand{\saegeblatt}[4]{%
  \begin{scope}[shift={(#1,#2)}]
    \pgfmathsetmacro{\sbs}{360/#4}%
    \pgfmathsetmacro{\sbz}{min(0.18*#3, 2*pi*#3/#4*0.55)}%
    \fill[black!60, draw=black!85, line width=0.4pt] (0:#3)
      \foreach \i in {1,...,#4}{
        -- ({\i*\sbs-\sbs}:#3) -- ({\i*\sbs-0.2*\sbs}:#3-\sbz) -- ({\i*\sbs}:#3)}
      -- cycle;
    \fill[black!25] (0,0) circle (#3*0.5);
    \fill[white] (0,0) circle (#3*0.12);
  \end{scope}}


% Nummer eines Bauteils wie in der Übersichtsskizze, z.B. \teil{3}
\newcommand{\teil}[1]{\tikz[baseline=(n.base)]\node[circle, draw=black, fill=white, line width=0.8pt, minimum size=3.8mm, inner sep=0pt, font=\tiny\bfseries](n){#1};}
% Sicherheitszeichen (ISO 7010) aus fablab-document
\newcommand{\zeichen}[2][1.4cm]{\includegraphics[height=#1]{\basedir/ba-symbole/#2}}
% Sägeblatt als kleines Symbol, z.B. \blattsymbol{36}
\newcommand{\blattsymbol}[1]{\tikz[x=1mm, y=1mm, baseline=-1mm]{\saegeblatt{0}{0}{5}{#1}}}
% Diamantblatt: geschlossener Rand mit Diamantbelag
\newcommand{\diamantsymbol}{\tikz[x=1mm, y=1mm, baseline=-1mm]{%
	\fill[black!60, draw=black!85, line width=0.4pt] (0,0) circle (5);
	\foreach \w in {0,15,...,345} \fill[white!80!gray] (\w:4.4) circle (0.35);
	\fill[black!25] (0,0) circle (2.5); \fill[white] (0,0) circle (0.6);}}

\date{\today}
% Versionsnummer: version.tex wird von der GitHub Action erzeugt,
% lokal steht stattdessen "Entwurf" mit Datum in der Fußzeile.
\InputIfFileExists{version}{}{\newcommand{\Version}{Entwurf \today}}
\fancyfoot[R]{Version \Version}
\author{kontakt@fablab.fau.de}
\title{Einweisung Tischkreissäge}

\begin{document}

	\maketitle
	\begin{center}
		Für die Feinschnitt-Tischkreissäge \textbf{Proxxon FET}\\[1mm]
		{\small Version \Version}
	\end{center}

	\textbf{Nur eingewiesene Personen dürfen die Tischkreissäge benutzen.}
	Wenn du noch nicht eingewiesen bist, frage einen Betreuer. Nach der Einweisung und Unterschrift auf der Einweisungsliste darfst du die Säge selbstständig verwenden.

	\begin{center}
		\zeichen{M003}\quad\zeichen{M004}\quad\zeichen{P028}\quad\zeichen{W024}
	\end{center}

	\section{Regeln und Sicherheit}

	\begin{itemize}
		\item Nie ohne \textbf{Spaltkeil} und \textbf{Sägeblattschutz} sägen. Beide dürfen nicht entfernt, verbogen oder festgeklemmt werden.
		\item Das Blatt nur so hoch stellen, dass es \textbf{knapp über das Werkstück} hinausragt.
		\item Hände nie in die Verlängerung des Sägeblatts bringen, weder vor noch hinter dem Blatt. Kleine und schmale Werkstücke mit dem \textbf{Schiebestock} schieben.
		\item Das Werkstück immer am Längsanschlag oder am Winkelanschlag führen. \textbf{Nie freihändig} sägen und \textbf{nie beide Anschläge gleichzeitig} verwenden.
		\item Seitlich neben der Schnittlinie stehen, nicht direkt dahinter (Rückschlag).
		\item Gehörschutz und Schutzbrille tragen. Beim Sägen \textbf{keine Handschuhe}, keine weite Kleidung, lange Haare zusammenbinden.
		\item Immer die \textbf{Absaugung} anschließen.
		\item Reste und Späne erst entfernen, wenn das Blatt steht. Nie in das auslaufende Blatt greifen.
		\item Vor allen Einstellarbeiten und vor dem Sägeblattwechsel den \textbf{Netzstecker ziehen}.
		\item Das Sägeblatt wechselt nur ein Betreuer.
		\item Schäden sofort einem Betreuer oder per Mail an \texttt{kontakt@fablab.fau.de} melden und die Säge nicht weiter benutzen.
	\end{itemize}

	\BAfalls{betriebsanweisung/ba_tischkreissaege}{Zusätzlich gilt die \textbf{Betriebsanweisung} auf der nächsten Seite (Abschnitt~\ref{sec:betriebsanweisung}). Sie hängt an der Säge aus und ist Bestandteil dieser Einweisung.}

	\newpage

	% Betriebsanweisung als ganze Seite im Format des Aushangs an der Säge
	\BAseite[sec:betriebsanweisung]{betriebsanweisung/ba_tischkreissaege}

	\renewcommand{\contentsname}{Inhaltsverzeichnis / Arbeitsablauf}
	\setcounter{tocdepth}{2}
	\tableofcontents
	\newpage

	\section{Die Säge und ihre Anleitung}

	\begin{center}
		\begin{tabular}{r|l}
			Hersteller & Proxxon \\
			Produktname & Feinschnitt-Tischkreissäge FET \\
			Leistung & $200\,\mathrm{W}$ \\
			Leerlaufdrehzahl & $7.000\,\mathrm{min}^{-1}$ \\
			Tischgröße & $300 \times 300\,\mathrm{mm}$, seitlich ausziehbar \\
			Schnitthöhe & $1 - 22\,\mathrm{mm}$ (bei geneigtem Blatt weniger) \\
			Neigung des Sägeblatts & $0$\textdegree{} bis $45$\textdegree, nach rechts \\
			Sägeblätter & $\varnothing\,50 - 85\,\mathrm{mm}$, Bohrung $10\,\mathrm{mm}$ \\
		\end{tabular}
	\end{center}

	Diese Einweisung fasst die Regeln des FabLabs zusammen. Sie ersetzt nicht die Betriebsanleitung des Herstellers: Die \textbf{Originalbetriebsanleitung} liegt der Maschine bei und ist auf der Produktseite von Proxxon zu finden.

	\begin{center}
		\qrcode[height=2.2cm]{https://www.proxxon.com/de/micromot/27070.php}\\[1mm]
		\href{https://www.proxxon.com/de/micromot/27070.php}{Proxxon FET bei Proxxon}
	\end{center}

	\textbf{Bestimmungsgemäße Verwendung:} Die FET ist eine kleine Tischkreissäge für genaue, gerade Schnitte in dünnen Werkstücken bis 22\,mm Dicke, etwa im Modellbau. Je nach Sägeblatt sägt sie Holz und Holzwerkstoffe, Kunststoffe, Aluminium und andere Nichteisenmetalle, Leiterplatten aus Epoxid (GFK) und Keramik. Sie wird dafür fest auf der Werkbank befestigt und immer mit Absaugung betrieben.

	\subsection{Die wichtigsten Teile}
	\label{sec:teile}

	\begin{center}
		% =====================================================================
%  Schematische Übersicht der Feinschnitt-Tischkreissäge Proxxon FET:
%  oben Draufsicht, unten Vorderansicht. Selbst gezeichnet (TikZ),
%  nicht maßstäblich; die Nummern zeigen, welche Teile es gibt, nicht
%  ihre exakte Lage an der Maschine. Nummern wie in der Tabelle der
%  Einweisung (\teil{n}). Lizenz: CC BY-SA 3.0
% =====================================================================
\begin{tikzpicture}[x=1mm, y=1mm, font=\sffamily\scriptsize]
  % ---------------------------------------------------------- Draufsicht
  \node[anchor=west, font=\small\bfseries] at (-40,100) {Draufsicht};
  % ausziehbarer Tisch links mit Hilfsanschlag
  \draw[metall] (-30,24) rectangle (6,26.5);
  \draw[metall] (-30,70) rectangle (6,72.5);
  \draw[gehaeuse] (-36,14) rectangle (-28,84);
  \draw[gelb] (-31,16) rectangle (-27,82);
  % Gehäuse und Sägetisch
  \draw[gehaeuse] (6,4) rectangle (92,90);
  \draw[tisch] (10,10) rectangle (88,86);
  % Führungsnuten für den Winkelanschlag
  \draw[black!60, line width=1.2pt] (20,10) -- (20,86) (66,10) -- (66,86);
  % Tischeinlage, Sägeblatt, Spaltkeil, Sägeblattschutz
  \draw[fill=black!70, draw=black] (41,26) rectangle (45,70);
  \fill[black!40] (42.4,32) rectangle (43.6,58);
  \draw[keil] (42.5,58) rectangle (43.5,68);
  \draw[draw=schutzorange!80!black, fill=schutzorange, fill opacity=0.55, line width=0.7pt]
    (40.5,30) rectangle (45.5,68);
  % Winkelanschlag in der linken Nut
  \begin{scope}[shift={(20,46)}]
    \draw[gelb] (-8,0) arc[start angle=180, end angle=0, radius=8] -- cycle;
    \draw[gelb] (-9,0) rectangle (15,2.5);
    \fill[black!70] (0,4) circle (1.6);
  \end{scope}
  % Längsanschlag rechts mit Klemmung und Feineinstellung
  \draw[metall] (72,4) rectangle (77,90);
  \draw[gelb] (73,-1) rectangle (76,4);
  \fill[black!80] (80,1) circle (2.2);
  \draw[black!80, line width=1pt] (80,1) -- (85,-2);
  % Skala vorn
  \fill[white] (10,4.5) rectangle (88,8.5);
  \foreach \x in {12,14,...,86} \draw[black!70, line width=0.3pt] (\x,8.5) -- (\x,7.2);
  \foreach \x in {12,22,...,82} \draw[black!70, line width=0.4pt] (\x,8.5) -- (\x,6.2);
  % Absauganschluss hinten
  \draw[metall] (43,94) circle (3.2);
  \fill[black!60] (43,94) circle (1.8);
  % Vorschub
  \draw[vorschub, blue!60!black] (52,14) -- (52,30);
  \node[anchor=west, blue!60!black] at (53,18) {Vorschub};
  % Nummern
  \nr{1}{30,62}{43,52}
  \nr{2}{56,77}{44.5,66}
  \nr{3}{82,80}{80,70}
  \nr{4}{100,60}{77,60}
  \nr{5}{8,52}{14,47}
  \nr{6}{60,40}{66,36}
  \nr{7}{-42,60}{-31,60}
  \nr{8}{30,0}{30,6.5}
  \nr{12}{56,98}{46,95}

  % ------------------------------------------------------ Vorderansicht
  \begin{scope}[shift={(0,-58)}]
    \node[anchor=west, font=\small\bfseries] at (-40,40) {Vorderansicht};
    % Tisch, Skala, Blatt und Anschlag über dem Tisch
    \draw[tisch] (6,30) rectangle (92,33);
    \draw[metall] (72,33) rectangle (77,39);
    \fill[black!60] (42.4,33) rectangle (43.6,37);
    \draw[draw=schutzorange!80!black, fill=schutzorange, fill opacity=0.55, line width=0.7pt]
      (40.5,33) rectangle (45.5,40);
    \draw[gelb] (-36,28) rectangle (-28,36);
    \draw[metall] (-28,31) rectangle (6,32.5);
    % Gehäuse
    \draw[gehaeuse] (6,0) rectangle (92,30);
    \fill[white] (10,25) rectangle (88,29);
    \foreach \x in {12,14,...,86} \draw[black!70, line width=0.3pt] (\x,29) -- (\x,27.7);
    % Schalter
    \draw[fill=black!80] (11,7) rectangle (21,21);
    \draw[fill=okgruen!70] (12.5,15) rectangle (19.5,19.5);
    \draw[fill=nokrot!80] (12.5,8.5) rectangle (19.5,13);
    % Fenster mit Winkelskala und Verstellung
    \draw[fill=gehaeusegruen!40, draw=black!80] (30,3) rectangle (68,23);
    \draw[gelb] ($(49,9)+(20:13)$) arc[start angle=20, end angle=160, radius=13]
      -- ($(49,9)+(160:10)$) arc[start angle=160, end angle=20, radius=10] -- cycle;
    \foreach \w in {25,35,...,155} \draw[black!70] ($(49,9)+(\w:13)$) -- ($(49,9)+(\w:11.6)$);
    \fill[black!85] (49,9) circle (4.2);
    \foreach \w in {0,60,...,300} \fill[black!85] ($(49,9)+(\w:4.2)$) circle (1.3);
    \fill[black!25] (49,9) circle (2);
    \draw[nokrot, line width=0.8pt] (49,9) -- ++(80:12.5);
    % Längsanschlag: Feineinstellung
    \fill[black!80] (80,32) circle (2);
    % Nummern
    \nr{9}{2,-6}{16,7}
    \nr{10}{40,-6}{48,7}
    \nr{11}{70,-6}{60.5,15}
  \end{scope}
  \node[font=\scriptsize] at (49,-70) {Bedienperson};
\end{tikzpicture}
\\[1mm]
		{\footnotesize Schematische Skizze, nicht maßstäblich. Die Lage einzelner Teile an der Maschine kann abweichen.}
	\end{center}

	\begin{center}
		\small
		\renewcommand{\arraystretch}{1.25}
		\renewcommand{\tabularxcolumn}[1]{m{#1}}
		\begin{tabularx}{\textwidth}{|c|m{42mm}|X|}
			\hline
			& \textbf{Teil} & \textbf{Wozu} \\ \hline
			\teil{1} & Sägeblatt & Ragt durch die Tischeinlage nach oben. Höhe und Neigung werden vorn am Gehäuse eingestellt \teil{10}. \\ \hline
			\teil{2} & Spaltkeil mit Sägeblattschutz & Der Spaltkeil hält den Schnitt hinter dem Blatt offen. Der Schutz darüber hebt sich beim Sägen selbst an und fällt danach zurück. Nie entfernen. \\ \hline
			\teil{3} & Sägetisch & Auflage für das Werkstück. Sauber und frei von Resten halten. \\ \hline
			\teil{4} & Längsanschlag & Führt das Werkstück bei Längsschnitten. Lässt sich von rechts oder links einsetzen und mit einer Stellschraube fein verstellen. \\ \hline
			\teil{5} & Winkelanschlag & Für Quer- und Gehrungsschnitte. Läuft in einer der Führungsnuten. \\ \hline
			\teil{6} & Führungsnuten & Links und rechts vom Blatt, für den Winkelanschlag. \\ \hline
			\teil{7} & Ausziehbarer Tisch mit Hilfsanschlag & Vergrößert die Auflage für breite Werkstücke. Die gelbe Kante dient dann als Anschlag. \\ \hline
			\teil{8} & Skala & Zum Einstellen des Längsanschlags. \\ \hline
			\teil{9} & Ein-/Ausschalter & Grün ein, rot aus. Im Notfall sofort auf Aus. \\ \hline
			\teil{10} & Verstellung für Höhe und Neigung & Klemmschraube lösen, einstellen, wieder festziehen (Abschnitt~\ref{sec:einstellungen}). \\ \hline
			\teil{11} & Winkelskala & Zeigt die Neigung des Sägeblatts. \\ \hline
			\teil{12} & Absauganschluss & Saugschlauch aufstecken. Nie ohne Absaugung sägen. \\ \hline
		\end{tabularx}
	\end{center}

	\section{Vorbereitung}

	\subsection{Checkliste vor dem Sägen}
	Vor \textbf{jedem} Einsatz durchgehen:
	\begin{itemize}
		\item Säge mit Zwingen fest auf der Werkbank? Tisch frei von Werkzeug und Resten?
		\item Netzkabel unbeschädigt und außerhalb des Schnittbereichs?
		\item Passendes Sägeblatt eingesetzt (Abschnitt~\ref{sec:saegeblatt}), scharf, ohne Risse und ohne fehlende Zähne?
		\item Spaltkeil und Sägeblattschutz montiert? Fällt der Schutz von selbst zurück?
		\item Absaugung angeschlossen, Sauger auf \enquote{AUTO}?
		\item Werkstück frei von Nägeln, Schrauben und anderen Metallteilen? Gerade Kante zum Anlegen vorhanden?
		\item Schnitthöhe und Neigung eingestellt, alle Klemmschrauben und Anschläge fest?
		\item Schiebestock griffbereit, Schutzausrüstung angelegt?
		\item Keine Unbeteiligten neben oder hinter der Säge?
	\end{itemize}

	\subsection{Schutzausrüstung}

	\begin{center}
		\begin{tabular}{cccc}
			\zeichen{M003} & \zeichen{M004} & \zeichen{P028} & \zeichen{M009} \\
			\small Gehörschutz & \small Schutzbrille & \small beim Sägen & \small nur beim \\
			& & \small keine Handschuhe & \small Blattwechsel
		\end{tabular}
	\end{center}

	\begin{itemize}
		\item \textbf{Gehörschutz:} hängt an der Fräse links neben der Werkbank.
		\item \textbf{Schutzbrille:} immer, besonders bei Aluminium, Faserwerkstoffen und Leiterplatten. Brillen hängen hinter der Werkbank und hinter dem Chemiearbeitsbereich an der Wand.
		\item \textbf{Staubmaske} bei viel Staub, z.\,B. bei MDF, Hartholz und Leiterplatten.
		\item Beim Sägen \textbf{keine Handschuhe} tragen: Das Blatt kann sie erfassen und die Hand mitziehen. \textbf{Schutzhandschuhe} (Schublade W8) nur bei gezogenem Netzstecker, z.\,B. beim Sägeblattwechsel oder beim Tragen von rauem Material.
		\item Eng anliegende Kleidung, keine Schals, Kordeln oder Schmuck; lange Haare zusammenbinden.
	\end{itemize}

	\subsection{Säge aufstellen und Absaugung anschließen}
	\begin{enumerate}
		\item Säge auf die Werkbank stellen und mit Zwingen festklemmen (Schublade W4). Genaues und sicheres Sägen geht nur, wenn die Säge nicht wackeln oder rutschen kann.
		\item Den Festool-Sauger an den Absauganschluss~\teil{12} anschließen.
		\item Das Netzkabel der Säge in die Steckdose am Sauger stecken und den Sauger auf \enquote{AUTO} stellen. Er läuft dann automatisch mit der Säge an und kurz nach.
		\item Vor dem Einstecken prüfen, dass der Schalter~\teil{9} auf Aus steht.
	\end{enumerate}
	Bei Holz können Späne den Absauganschluss zusetzen. Lässt die Absaugung nach, die Säge ausschalten, Netzstecker ziehen und die Verstopfung beseitigen.

	\subsection{Material}
	\begin{itemize}
		\item Zugelassen: Holz und Holzwerkstoffe, Kunststoffe, Aluminium und andere weiche Nichteisenmetalle, Leiterplatten und Keramik, jeweils mit dem passenden Blatt.
		\item Höchstens 22\,mm dick, bei geneigtem Blatt weniger.
		\item Nicht zugelassen: Stahl und Eisen, Material mit Nägeln, Schrauben oder Klammern.
		\item Keine Werkstücke, die sich nicht sicher am Anschlag führen lassen (sehr klein, rund, stark verzogen). Im Zweifel einen Betreuer fragen.
		\item Der Staub mancher Materialien ist gesundheitsschädlich, vor allem von Hartholz, MDF und GFK (Leiterplatten). Bei Leiterplatten muss die Absaugung unbedingt laufen.
	\end{itemize}

	\subsection{Sägeblatt wählen}
	\label{sec:saegeblatt}
	Das Blatt muss zum Material passen. Viele Zähne schneiden langsamer, aber sauberer; wenige Zähne schneiden schneller und gröber. Für die Säge gibt es im FabLab diese Blätter:

	\begin{center}
		\small
		\renewcommand{\arraystretch}{1.4}
		\renewcommand{\tabularxcolumn}[1]{m{#1}}
		\begin{tabularx}{\textwidth}{|c|m{38mm}|X|}
			\hline
			& \textbf{Sägeblatt} & \textbf{Material} \\ \hline
			\blattsymbol{80} & Super-Cut, sehr viele feine Zähne & Weiche Hölzer und Kunststoffe, wenn der Schnitt besonders sauber werden soll \\ \hline
			\blattsymbol{36} & Hartmetall, 36 Zähne & Balsa, Sperrholz, Weich- und Hartholz, Kunststoff, Aluminium \\ \hline
			\blattsymbol{24} & Hartmetall, 24 Zähne & Hartholz, Spanplatte, Kunststoff, Aluminium \\ \hline
			\diamantsymbol & Diamantblatt, ohne Zähne & Keramik und Leiterplatten (GFK) \\ \hline
		\end{tabularx}
	\end{center}

	Passt das eingesetzte Blatt nicht zum Material, einen Betreuer fragen. Den Wechsel macht ein Betreuer (Abschnitt~\ref{sec:blattwechsel}).

	\section{Einstellungen}
	\label{sec:einstellungen}

	Vor allen Einstellarbeiten die Säge ausschalten und den \textbf{Netzstecker ziehen}.

	\subsection{Schnitthöhe}
	Das Blatt soll nur \textbf{knapp über das Werkstück} hinausragen, etwa eine Zahnhöhe. Dann liegt wenig Blatt frei, und der Schnitt wird sauberer.

	\begin{center}
		% =====================================================================
%  Schnitthöhe: Das Blatt ragt nur knapp über das Werkstück hinaus.
%  Seitenansicht, Bedienperson rechts, Vorschub nach links.
%  Selbst gezeichnet (TikZ), Lizenz: CC BY-SA 3.0
% =====================================================================
\begin{tikzpicture}[x=1mm, y=1mm, font=\scriptsize]

  % Seitenansicht mit Werkstück der Dicke 8: \saegeseite{Blattüberstand}
  % (Tischoberkante bei y=14, Blattmitte bei x=40)
  \newcommand{\saegeseite}[1]{%
    \pgfmathsetmacro{\mitte}{14+8+#1-17}%
    \begin{scope}
      \clip (0,14) rectangle (82,46);
      \saegeblatt{40}{\mitte}{17}{24}
    \end{scope}
    % Spaltkeil: folgt dem Blatt im Abstand von etwa 2 mm
    \pgfmathsetmacro{\wa}{180-asin((14-\mitte)/19)}%
    \draw[keil] ($(40,\mitte)+(\wa:19)$) arc[start angle=\wa, end angle=112, radius=19]
      -- ($(40,\mitte)+(112:21.5)$) arc[start angle=112, end angle=\wa, radius=21.5] -- cycle;
    \fill[gehaeusegruen!70] (4,4) rectangle (78,12);
    \draw[tisch] (4,12) rectangle (78,14);
    \draw[werkstueck, fill opacity=0.6] (8,14) rectangle (76,22);
    \draw[vorschub] (70,26) -- (56,26);}

  % ---------------- 1 richtig --------------------------------------
  \feld{0}{0}{82}{52}{1\quad Blatt knapp über dem Werkstück}
  \saegeseite{3}
  \draw[mass] (60,22) -- (60,25) node[right=-0.5mm, midway] {};
  \draw[hilfslinie] (44,25) -- (62,25);
  \node[anchor=west, align=left] at (62,33) {knapp eine\\Zahnhöhe};
  \draw[hilfslinie] (61.5,30.5) -- (60,24);
  \pic at (74,45) {ok};
  \node at (41,7.5) {\color{white}wenig Blatt frei, sauberer Schnitt};

  % ---------------- 2 falsch ---------------------------------------
  \begin{scope}[shift={(88,0)}]
    \feld{0}{0}{82}{52}{2\quad Blatt zu hoch}
    \saegeseite{15}
    \pic at (74,45) {nok};
    \node at (41,7.5) {\color{white}viel freies Blatt, höhere Verletzungsgefahr};
  \end{scope}
\end{tikzpicture}

	\end{center}

	\begin{enumerate}
		\item Die große Klemmschraube vorn am Gehäuse~\teil{10} lösen.
		\item Mit der kleinen Stellschraube die Höhe einstellen: im Uhrzeigersinn fährt das Blatt nach oben, gegen den Uhrzeigersinn nach unten.
		\item Werkstück neben das Blatt halten und die Höhe prüfen.
		\item Klemmschraube wieder festziehen.
	\end{enumerate}

	\subsection{Neigung für Gehrungsschnitte}
	Das Blatt lässt sich für Gehrungsschnitte von 0\textdegree{} bis 45\textdegree{} nach rechts neigen.
	\begin{enumerate}
		\item Klemmung der Neigung~\teil{10} lösen.
		\item Blatt nach rechts schwenken, bis der Zeiger an der Winkelskala~\teil{11} den gewünschten Winkel zeigt.
		\item Klemmung wieder festziehen.
		\item Den Längsanschlag auf die \textbf{linke} Seite setzen, damit sich das Blatt vom Anschlag weg neigt.
	\end{enumerate}

	\begin{center}
		% =====================================================================
%  Geneigtes Sägeblatt (Vorderansicht, Blick der Bedienperson): Das Blatt
%  neigt sich nach rechts. Den Längsanschlag dann links vom Blatt setzen,
%  damit kein Reststück unter dem Blatt eingeklemmt wird.
%  Selbst gezeichnet (TikZ), Lizenz: CC BY-SA 3.0
% =====================================================================
\begin{tikzpicture}[x=1mm, y=1mm, font=\scriptsize]

  % Tisch mit um 30° geneigtem Blatt, Drehpunkt auf der Tischoberfläche
  \newcommand{\neigetisch}{%
    \fill[gehaeusegruen!70] (6,4) rectangle (76,12);
    \draw[tisch] (6,12) rectangle (76,14);
    \draw[hilfslinie, dashed] (41,14) -- (41,36);}
  \newcommand{\neigeblatt}{%
    \draw[blatt, rotate around={-30:(41,14)}] (40.4,14) rectangle (41.6,34);}

  % ---------------- 1 richtig --------------------------------------
  \feld{0}{0}{82}{52}{1\quad Anschlag links vom Blatt}
  \neigetisch
  \draw[werkstueck] (18,14) -- (41,14) -- (45.6,22) -- (18,22) -- cycle;
  \draw[werkstueck] (41.9,14) -- (60,14) -- (60,22) -- (46.5,22) -- cycle;
  \neigeblatt
  \draw[metall] (12,14) rectangle (18,26);
  \draw[drehung] (41,30) arc[start angle=90, end angle=60, radius=16];
  \node at (50,37) {0--45°};
  \node[anchor=west, align=left] at (60,30) {Blatt neigt\\sich vom\\Anschlag weg};
  \node[anchor=south] at (15,26) {Anschlag};
  \pic at (74,45) {ok};

  % ---------------- 2 falsch ---------------------------------------
  \begin{scope}[shift={(88,0)}]
    \feld{0}{0}{82}{52}{2\quad Anschlag rechts vom Blatt}
    \neigetisch
    \draw[werkstueck] (14,14) -- (41,14) -- (45.6,22) -- (14,22) -- cycle;
    % Reststück unter dem geneigten Blatt eingeklemmt
    \draw[werkstueck] (41,14) -- (51,14) -- (51,22) -- (45.6,22) -- cycle;
    \draw[metall] (51,14) rectangle (57,26);
    \neigeblatt
    \draw[kraft, nokrot] (48,24) -- (48,34);
    \node[nokrot, anchor=west, align=left] at (58,32) {Reststück\\klemmt unter\\dem Blatt};
    \pic at (74,45) {nok};
  \end{scope}
\end{tikzpicture}

	\end{center}

	Die Skala ist nur ungefähr genau. Für genaue Winkel einen Probeschnitt in einem Reststück machen und mit dem Winkelmesser nachmessen.

	\subsection{Längsanschlag}
	Am Längsanschlag~\teil{4} wird das Werkstück entlanggeschoben. Der Abstand zwischen Blatt und Anschlag ist die Breite des abgesägten Streifens.
	\begin{enumerate}
		\item Anschlag von rechts oder links in die Führung am Tisch schieben. Zum Einsetzen, Verschieben und Herausnehmen beide Klemmungen lösen.
		\item Anschlag ungefähr auf das Maß schieben, an der Skala~\teil{8} ablesen. Für genaue Maße mit der Stellschraube nachstellen.
		\item Beide Klemmungen festziehen. Der Anschlag darf sich beim Sägen nicht bewegen.
		\item Das Maß am besten direkt zwischen Anschlag und einem Zahn des Blatts nachmessen.
	\end{enumerate}

	\subsection{Ausziehbarer Tisch und Hilfsanschlag}
	Für breitere Werkstücke lässt sich der Tisch zur Seite ausziehen~\teil{7}.
	\begin{enumerate}
		\item Die gelbe Kante abnehmen.
		\item Den Tisch so weit wie nötig herausziehen und mit der Stütze abstützen.
		\item Mit der kleinen Klemmschraube feststellen.
		\item Für eine ebene Auflage die gelbe Kante bündig einsetzen. Soll sie als \textbf{Hilfsanschlag} dienen, bleibt sie hochgestellt: Ihr Abstand zum Blatt ist dann die Schnittbreite. Vor dem Sägen fest klemmen.
	\end{enumerate}

	\subsection{Winkelanschlag}
	Für Querschnitte und schräge Schnitte wird das Werkstück am Winkelanschlag~\teil{5} geführt.
	\begin{enumerate}
		\item Winkelanschlag in die linke oder rechte Führungsnut~\teil{6} einsetzen.
		\item Klemmknopf lösen, Winkel an der Skala einstellen, Klemmknopf festziehen.
		\item Den Längsanschlag abnehmen oder weit vom Blatt wegschieben.
	\end{enumerate}

	\section{Sägen}

	\begin{center}
		% =====================================================================
%  Werkstück führen: Längsschnitt mit Längsanschlag und Schiebestock,
%  Querschnitt mit Winkelanschlag und zwei typische Fehler.
%  Draufsicht, Bedienperson unten, Vorschub nach oben.
%  Selbst gezeichnet (TikZ), Lizenz: CC BY-SA 3.0
% =====================================================================
\begin{tikzpicture}[x=1mm, y=1mm, font=\scriptsize]

  % Hand von oben: \hand{x}{y}{Drehwinkel}
  \newcommand{\hand}[3]{%
    \begin{scope}[shift={(#1,#2)}, rotate=#3, scale=0.75]
      \foreach \fx/\fl in {-3.4/4.2, -1.3/5.2, 0.8/5, 2.8/4}
        \draw[hand, rounded corners=0.8mm] (\fx-0.9,1) rectangle (\fx+0.9,1+\fl+2);
      \draw[hand, rounded corners=0.8mm, rotate around={40:(-4,-1)}] (-5,-1) rectangle (-3.2,5);
      \draw[hand] (-4.4,-5) rectangle (3.8,2.6);
    \end{scope}}
  % Tisch mit Sägeblatt (Draufsicht): \tischblatt
  \newcommand{\tischblatt}{%
    \draw[tisch] (8,4) rectangle (74,44);
    \draw[black!60, line width=1pt] (20,4) -- (20,44);
    \fill[black!70] (39.5,16) rectangle (42.5,40);
    \fill[black!40] (40.4,20) rectangle (41.6,30);
    \draw[keil] (40.5,31) rectangle (41.5,37);}
  % Gefahrenzone vor und hinter dem Blatt
  \newcommand{\zone}{\fill[gefahrenzone] (37,4) rectangle (45,44);}

  % ---------------- 1 Längsschnitt ---------------------------------
  \feld{0}{57}{82}{52}{1\quad Längsschnitt}
  \begin{scope}[shift={(0,57)}]
    \tischblatt
    \zone
    \draw[metall] (54,4) rectangle (57,44);
    % Werkstück zwischen Blatt und Anschlag, schon zur Hälfte gesägt
    \draw[werkstueck] (26,22) -- (54,22) -- (54,44) -- (41.6,44) -- (41.6,30)
      -- (40.4,30) -- (40.4,44) -- (26,44) -- cycle;
    % Schiebestock schiebt den Streifen
    \draw[fill=holzrand!70, draw=holzrand!60!black, line width=0.5pt]
      (43,22) -- (49,22) -- (49,19) -- (60,7) -- (57,4.5) -- (45.5,17.5) -- (43,19.5) -- cycle;
    \hand{60}{5}{-40}
    \draw[vorschub] (66,20) -- (66,32);
    \node[anchor=west, align=left] at (58,40) {Längs-\\anschlag};
    \node[anchor=east, align=right] at (36,12) {Schiebestock\\für die letzten cm};
    \draw[hilfslinie] (36.5,12) -- (50,13);
    \pic at (74,47) {ok};
  \end{scope}

  % ---------------- 2 Hand vor dem Blatt ---------------------------
  \feld{88}{57}{82}{52}{2\quad Freihändig, Hand vor dem Blatt}
  \begin{scope}[shift={(88,57)}]
    \tischblatt
    \zone
    \draw[werkstueck, rotate around={8:(41,26)}] (26,18) rectangle (56,34);
    \hand{42}{12}{0}
    \node[nokrot, anchor=west, align=left, fill=white, fill opacity=0.85, text opacity=1, inner sep=1pt]
      at (50,10) {Hand in der\\Blattverlängerung};
    \pic at (74,47) {nok};
  \end{scope}

  % ---------------- 3 Querschnitt mit Winkelanschlag -----------------
  \feld{0}{0}{82}{52}{3\quad Querschnitt mit Winkelanschlag}
  \begin{scope}[shift={(0,0)}]
    \tischblatt
    % Längsanschlag weit weg
    \draw[metall, opacity=0.35, dashed] (66,4) rectangle (69,44);
    % Winkelanschlag in der linken Nut, Werkstück davor
    \begin{scope}[shift={(20,14)}]
      \draw[gelb] (-7,0) arc[start angle=180, end angle=0, radius=7] -- cycle;
      \draw[gelb] (-8,0) rectangle (16,-2.5);
      \fill[black!70] (0,3.5) circle (1.4);
    \end{scope}
    \draw[werkstueck] (13,14.2) rectangle (54,22);
    \hand{16}{6}{0}
    \draw[vorschub] (60,8) -- (60,20);
    \node[anchor=west, align=left] at (52,31) {Längsanschlag\\weit weg oder\\abgenommen};
    \pic at (74,47) {ok};
  \end{scope}

  % ---------------- 4 beide Anschläge -----------------------------
  \feld{88}{0}{82}{52}{4\quad Beide Anschläge zugleich}
  \begin{scope}[shift={(88,0)}]
    \tischblatt
    \draw[metall] (52,4) rectangle (55,44);
    \begin{scope}[shift={(20,14)}]
      \draw[gelb] (-7,0) arc[start angle=180, end angle=0, radius=7] -- cycle;
      \draw[gelb] (-8,0) rectangle (16,-2.5);
      \fill[black!70] (0,3.5) circle (1.4);
    \end{scope}
    \draw[werkstueck] (13,14.2) rectangle (40.4,22);
    % Reststück verkeilt sich zwischen Blatt und Anschlag
    \draw[werkstueck, rotate around={-25:(47,24)}] (41.6,20) rectangle (52,28);
    \draw[kraft, nokrot] (47,30) -- (47,40) node[above, align=center] {};
    \node[nokrot, anchor=west, align=left, fill=white, fill opacity=0.85, text opacity=1, inner sep=1pt]
      at (57,24) {Reststück\\klemmt\\und fliegt};
    \pic at (74,47) {nok};
  \end{scope}
\end{tikzpicture}

	\end{center}

	\subsection{Längsschnitt mit Längsanschlag}
	\begin{enumerate}
		\item Checkliste durchgehen, Schutzausrüstung anlegen.
		\item Längsanschlag auf das Maß stellen und klemmen, Blatthöhe einstellen.
		\item Werkstück mit der geraden Kante an den Anschlag legen, noch ohne das Blatt zu berühren.
		\item Seitlich neben der Schnittlinie stehen, Säge einschalten~\teil{9}. Warten, bis das Blatt die volle Drehzahl hat.
		\item Werkstück langsam und gleichmäßig nach vorn schieben: kräftig auf den Tisch und leicht an den Anschlag drücken, nur wenig gegen das Blatt. Je dicker das Werkstück und je feiner die Zähne, desto langsamer.
		\item Sobald die Hand näher als eine Handbreite an das Blatt käme, mit dem \textbf{Schiebestock} weiterschieben. Schmale Streifen immer mit dem Schiebestock.
		\item Werkstück ganz durchschieben, bis es hinter dem Spaltkeil frei ist. Nicht zurückziehen, solange das Blatt läuft.
		\item Säge ausschalten und warten, bis das Blatt steht. Erst dann Reste entfernen.
	\end{enumerate}

	\subsection{Querschnitt mit Winkelanschlag}
	\begin{enumerate}
		\item Längsanschlag abnehmen oder weit wegschieben, damit das abgesägte Stück nicht eingeklemmt wird.
		\item Winkel am Winkelanschlag einstellen und klemmen.
		\item Werkstück an den Winkelanschlag legen und mit der Hand festhalten, die Hand weit weg von der Schnittlinie.
		\item Säge einschalten, Winkelanschlag mit dem Werkstück gleichmäßig am Blatt vorbeischieben.
		\item Säge ausschalten, Anschlag erst zurückziehen, wenn das Blatt steht.
	\end{enumerate}

	\subsection{Rückschlag}
	Ein Rückschlag entsteht, wenn das Werkstück am Blatt klemmt oder sich verdreht. Die hinteren Zähne steigen nach oben, heben das Werkstück an und schleudern es mit großer Kraft in Richtung der Bedienperson. Auch bei einer kleinen Säge wie der FET kann das Teil gefährlich schnell werden, und die Hand wird dabei leicht in das Blatt gezogen.

	\begin{center}
		% =====================================================================
%  Rückschlag: wie er entsteht, wozu der Spaltkeil da ist und wo man
%  steht. Feld 1 Seitenansicht (Bedienperson rechts, Vorschub nach
%  links), Felder 2-4 Draufsicht (Bedienperson unten, Vorschub nach oben).
%  Selbst gezeichnet (TikZ), Lizenz: CC BY-SA 3.0
% =====================================================================
\begin{tikzpicture}[x=1mm, y=1mm, font=\scriptsize]

  % ---------------- 1 Entstehung (Seitenansicht) -------------------
  \feld{0}{57}{82}{52}{1\quad So entsteht ein Rückschlag}
  \begin{scope}[shift={(0,57)}]
    \begin{scope}
      \clip (0,14) rectangle (82,46);
      \saegeblatt{40}{8}{17}{24}
    \end{scope}
    \fill[gehaeusegruen!70] (4,4) rectangle (78,12);
    \draw[tisch] (4,12) rectangle (78,14);
    % Werkstück wird hinten von den aufsteigenden Zähnen angehoben
    \draw[werkstueck, fill opacity=0.65, rotate around={-13:(64,14)}] (16,14) rectangle (64,21);
    % Drehrichtung des Blatts
    \draw[drehung, nokrot] ($(40,8)+(160:22)$) arc[start angle=160, end angle=120, radius=22];
    \node[nokrot, align=center] at (13,40) {Zähne hinten\\steigen auf};
    % Werkstück fliegt zur Bedienperson
    \draw[kraft, nokrot, line width=1.6pt] (60,30) -- (78,30);
    \node[nokrot, align=center] at (66,38) {Werkstück fliegt\\zur Bedienperson};
    \node[white] at (41,7.5) {Werkstück klemmt oder verdreht sich};
  \end{scope}

  % Draufsicht auf Blatt und Schnitt: \draufblatt{x}{y}
  \newcommand{\draufblatt}[2]{%
    \fill[black!60] (#1-0.6,#2-12) rectangle (#1+0.6,#2+12);
    \foreach \t in {-11,-9,...,11} \fill[black!80] (#1-0.9,#2+\t) rectangle (#1+0.9,#2+\t+0.8);}

  % ---------------- 2 ohne Spaltkeil -------------------------------
  \feld{88}{57}{82}{52}{2\quad Ohne Spaltkeil}
  \begin{scope}[shift={(88,57)}]
    % Schnittfuge schließt sich hinter dem Blatt und klemmt es ein
    \draw[werkstueck] (8,4) -- (74,4) -- (74,44) -- (41.1,44) -- (41.1,36) -- (41.8,24)
      -- (41.8,12) -- (40.2,12) -- (40.2,24) -- (40.9,36) -- (40.9,44) -- (8,44) -- cycle;
    \draw[vorschub] (60,8) -- (60,20);
    \draw[kraft, nokrot] (30,34) -- (39,34);
    \draw[kraft, nokrot] (52,34) -- (43,34);
    \node[nokrot, align=center, fill=white, fill opacity=0.8, text opacity=1, inner sep=1pt] at (23,24)
      {Schnitt schließt sich\\und klemmt das Blatt};
    \fill[black!60] (40.4,12) rectangle (41.6,26);
    \pic at (74,47) {nok};
  \end{scope}

  % ---------------- 3 mit Spaltkeil -------------------------------
  \feld{0}{0}{82}{52}{3\quad Mit Spaltkeil}
  \begin{scope}[shift={(0,0)}]
    \draw[werkstueck] (8,4) -- (74,4) -- (74,44) -- (41.8,44) -- (41.8,12) -- (40.2,12)
      -- (40.2,44) -- (8,44) -- cycle;
    \draw[vorschub] (60,8) -- (60,20);
    \fill[black!60] (40.4,12) rectangle (41.6,26);
    \draw[keil] (40.3,29) rectangle (41.7,40);
    \draw[hilfslinie] (41.7,35) -- (52,40);
    \node[anchor=west, align=left, fill=white, fill opacity=0.8, text opacity=1, inner sep=1pt] at (52,40)
      {Spaltkeil hält\\den Schnitt offen};
    \draw[hilfslinie] (40.4,19) -- (30,24);
    \node[anchor=east, fill=white, fill opacity=0.8, text opacity=1, inner sep=1pt] at (30,24) {Sägeblatt};
    \pic at (74,47) {ok};
  \end{scope}

  % ---------------- 4 Standposition -------------------------------
  \feld{88}{0}{82}{52}{4\quad Seitlich vom Blatt stehen}
  \begin{scope}[shift={(88,0)}]
    \draw[tisch] (14,14) rectangle (70,44);
    \fill[black!60] (40.4,26) rectangle (41.6,38);
    \draw[metall] (54,14) rectangle (57,44);
    % Rückschlagzone in der Verlängerung des Blatts
    \fill[gefahrenzone] (35,1) rectangle (47,26);
    % Bedienperson links neben der Zone
    \fill[black!45] (24,4) ellipse (8 and 3);
    \fill[black!65] (24,4.8) circle (2.6);
    \node[anchor=west, nokrot, align=left] at (50,6) {Rückschlagzone:\\hier nicht stehen};
    \draw[hilfslinie] (49.5,6) -- (47,8);
    \pic at (74,47) {ok};
  \end{scope}
\end{tikzpicture}

	\end{center}

	So wird er vermieden:
	\begin{itemize}
		\item Spaltkeil immer montiert lassen. Die Schnittfuge des Blatts muss mindestens so breit sein wie der Spaltkeil dick ist, sonst bleibt das Werkstück am Keil hängen.
		\item Werkstück immer an einem Anschlag führen, nie freihändig und nie gleichzeitig an Längs- und Winkelanschlag.
		\item Bei geneigtem Blatt den Längsanschlag links vom Blatt setzen.
		\item Seitlich neben der Schnittlinie stehen, nicht dahinter.
		\item Werkstück ohne Pause durchschieben und nicht während des Schnitts zurückziehen. Wenn es klemmt: ausschalten, Blatt stehen lassen, dann das Werkstück lösen.
		\item Nur scharfe Blätter verwenden. Ein stumpfes Blatt braucht mehr Druck und klemmt leichter.
	\end{itemize}

	\subsection{Kleine Teile und Leiterplatten}
	\begin{itemize}
		\item Kleine Teile nie mit den Fingern am Blatt vorbeischieben, sondern mit dem Schiebestock oder einem Hilfsbrett am Winkelanschlag.
		\item Kurze Abschnitte von einem langen Stück absägen, nicht ein kurzes Stück noch einmal teilen.
		\item Leiterplatten nur mit dem Diamantblatt und laufender Absaugung sägen, Staubmaske tragen. Der Glasfaserstaub reizt Haut, Augen und Atemwege.
	\end{itemize}

	\subsection{Nicht erlaubt}
	Folgendes ist im FabLab \textbf{verboten}:
	\begin{itemize}
		\item Sägen ohne Spaltkeil oder mit entferntem, manipuliertem oder festgeklemmtem Sägeblattschutz.
		\item Freihändiges Sägen ohne Anschlag.
		\item Gleichzeitiges Verwenden von Längs- und Winkelanschlag.
		\item Nuten, Falze und andere verdeckte Schnitte ohne Rücksprache mit einem Betreuer.
		\item Sägen mit Handschuhen.
		\item Sägen ohne Absaugung.
		\item Sägen von Stahl, Eisen oder Material mit Nägeln und Schrauben.
		\item Verwenden nicht passender, stumpfer, verbogener oder rissiger Sägeblätter.
		\item Benutzen der Säge durch nicht eingewiesene Personen.
	\end{itemize}

	\section{Nach dem Sägen}

	\begin{itemize}
		\item Säge ausschalten, Blatt abwarten, Netzstecker ziehen.
		\item Säge, Tisch und Arbeitsplatz gründlich aussaugen. \textbf{Nicht mit Druckluft abblasen}, das verteilt den Staub im Raum und drückt ihn in die Mechanik.
		\item Anschläge, Schiebestock, Zwingen und Sauger an ihren Platz zurücklegen.
		\item Schäden, stumpfe Sägeblätter oder Probleme einem Betreuer melden.
	\end{itemize}

	\section{Infos für Betreuer}

	\subsection{Sägeblattwechsel}
	\label{sec:blattwechsel}
	Nur durch Betreuer und nur mit gezogenem Netzstecker. Schutzhandschuhe tragen, die Zähne sind auch an alten Blättern sehr scharf. Vor dem Einschalten die Handschuhe wieder ausziehen. Die Innensechskantschlüssel liegen in der Maschinenkiste.
	\begin{enumerate}
		\item \textbf{Netzstecker ziehen.}
		\item Blatt ganz nach unten fahren (Abschnitt~\ref{sec:einstellungen}).
		\item Das Oberteil der Säge aufklappen.
		\item Welle blockieren: den kleinen Innensechskantschlüssel durch die Bohrung im Tisch in die Querbohrung der Sägewelle stecken. Dazu das Blatt von Hand drehen, bis die Bohrungen übereinander liegen.
		\item Mit dem großen Innensechskantschlüssel die Schraube in der Mitte des Blatts lösen und herausdrehen, das alte Blatt abnehmen.
		\item Neues Blatt aufsetzen. Die Bohrung muss sauber auf dem Bund der Welle sitzen, der Drehrichtungspfeil auf dem Blatt muss zur Drehrichtung der Säge passen.
		\item Schraube eindrehen und festziehen, solange die Welle noch blockiert ist.
		\item \textbf{Schlüssel aus der Welle ziehen}, Oberteil zuklappen und verriegeln.
		\item Prüfen: Blatt dreht sich von Hand frei, Spaltkeil steht in einer Linie mit dem Blatt, Sägeblattschutz fällt von selbst zurück.
	\end{enumerate}
	Keine deformierten, rissigen oder stumpfen Blätter einsetzen. Nur Blätter mit 50 bis 85\,mm Durchmesser und 10\,mm Bohrung, die laut Betriebsanleitung für die FET zugelassen sind.

	\subsection{Typische Fehler}
	\begin{itemize}
		\item \textbf{Brandspuren am Holz oder geschmolzener Kunststoff:} Vorschub zu langsam oder Blatt stumpf bzw.\ ungeeignet.
		\item \textbf{Ausgefranste Kanten:} Blatt mit mehr Zähnen nehmen, Blatthöhe prüfen.
		\item \textbf{Schnitt läuft schräg oder Blatt klemmt:} Anschlag prüfen und neu klemmen, Spaltkeil auf Ausrichtung prüfen.
		\item \textbf{Schwache Absaugung:} Absauganschluss und Schlauch auf Späne prüfen.
	\end{itemize}

	\subsection{Pflege und Prüfung}
	\begin{itemize}
		\item Säge, Tischeinlage und Absaugung regelmäßig von Staub und Spänen reinigen.
		\item Spaltkeil, Sägeblattschutz und Sägeblätter regelmäßig prüfen, stumpfe Blätter aussortieren.
		\item Netzkabel prüfen. Elektrische Prüfung nach DGUV Vorschrift 3 gemäß Prüffrist.
	\end{itemize}

	\section{Quellen}
	Die Formulierungen dieser Einweisung und alle Zeichnungen sind eigenständig erstellt (TikZ, Quelltext in \texttt{zeichnungen/}). Die Sicherheitszeichen nach ISO 7010 stammen aus \href{https://github.com/fau-fablab/fablab-document}{fablab-document} (gemeinfrei bzw.\ CC0, Wikimedia Commons). Technische Daten und sachliche Angaben beruhen auf öffentlich zugänglichen Herstellerangaben und allgemein bekannten Regeln für Tischkreissägen. Maßgeblich ist die Originalbetriebsanleitung von Proxxon, die der Maschine beiliegt (Link siehe oben). \enquote{Proxxon} ist eine Marke ihres Inhabers und wird hier nur zur Bezeichnung des Produkts verwendet. Dieses Dokument ist inoffiziell und wurde von Proxxon weder geprüft, noch übernimmt Proxxon dafür eine Haftung.

	\ccLicense{tischkreissaege-einweisung}{Einweisung Tischkreissäge}

\end{document}
